\chapter{Lung Cancer Genetic Markers}

\section*{What Is This Test?}
Lung cancer genetic marker testing (also known as molecular profiling or biomarker testing of non-small cell lung cancer [NSCLC]) detects somatic genomic alterations in tumor DNA that predict response or resistance to targeted therapy. The testing is performed on tumor tissue (biopsy or surgical resection) or on circulating tumor DNA (ctDNA, ``liquid biopsy''). The National Comprehensive Cancer Network (NCCN) \cite{nccn2024non} and the College of American Pathologists (CAP)/International Association for the Study of Lung Cancer (IASLC)/Association for Molecular Pathology (AMP) guidelines \cite{lindeman2018updated} recommend that all patients with advanced or metastatic non-small cell lung cancer (particularly adenocarcinoma and NSCLC not otherwise specified [NOS]) undergo comprehensive molecular profiling for the following actionable biomarkers: EGFR (epidermal growth factor receptor) mutations (exon 19 deletions, exon 21 L858R, exon 20 T790M, exon 20 insertions, and rare sensitizing mutations), ALK (anaplastic lymphoma kinase) rearrangements, ROS1 rearrangements, BRAF V600E mutation, NTRK1/2/3 fusions, RET rearrangements, MET exon 14 skipping mutations, KRAS G12C mutation, HER2 (ERBB2) mutations, and PD-L1 immunohistochemistry (for immunotherapy selection). In squamous cell carcinoma, fewer genetic alterations are actionable; however, EGFR, ALK, ROS1 testing is still recommended in never-smokers or light smokers, and comprehensive profiling is increasingly recommended. Testing is performed by next-generation sequencing (NGS) of a targeted gene panel, with reflex FISH or IHC for specific alterations if not covered by the NGS panel.

\section*{Alternative Names}
NSCLC Molecular Profiling; Lung Cancer Biomarker Testing; Lung Cancer Genomic Panel; EGFR/ALK/ROS1 Testing; Comprehensive Genomic Profiling in NSCLC; Liquid Biopsy for Lung Cancer

\section*{Principle}
NGS-based targeted panels simultaneously sequence DNA (for mutations and copy number changes) and RNA (for gene fusions including ALK, ROS1, RET, NTRK) from formalin-fixed paraffin-embedded tumor tissue. DNA sequencing detects single nucleotide variants, small insertions/deletions, and copy number changes in an actionable gene panel (EGFR, KRAS, BRAF, HER2, MET, PIK3CA, and others). RNA sequencing or anchored multiplex PCR detects actionable gene fusions (ALK, ROS1, RET, NTRK1/2/3) and MET exon 14 skipping. FISH using break-apart probes (ALK, ROS1, RET) provides an alternative method. IHC for ALK (D5F3 clone, Ventana) and ROS1 (D4D6 clone) serves as screening tests. PD-L1 IHC (22C3, 28-8, SP142, or SP263 antibody clones) assesses PD-L1 expression on tumor cells as a predictive biomarker for immunotherapy. EGFR ctDNA testing (cobas EGFR Mutation Test v2, Guardant360, FoundationOne Liquid CDx) can detect EGFR mutations from plasma when tissue is insufficient or unavailable.

\section*{History}
EGFR mutations in NSCLC were discovered in 2004 (Lynch et al., Paez et al., Pao et al.) in the EGFR tyrosine kinase domain (exons 18--21) as a predictor of dramatic response to gefitinib and erlotinib \cite{mok2009gefitinib}. ALK rearrangements (EML4-ALK fusion) in NSCLC were discovered by Soda et al. in 2007. ROS1 rearrangements were described in 2007--2009. Crizotinib, the first ALK/ROS1/MET inhibitor, was approved in 2011. The NCCN guidelines have progressively expanded the list of required biomarkers from EGFR alone (2011) to EGFR, ALK, and ROS1 (2012) to nine mandatory biomarkers (2023). The 2018 FDA approval of larotrectinib for NTRK fusions (tumor-agnostic) was followed by entrectinib for NTRK and ROS1 fusions. Sotorasib, the first KRAS G12C inhibitor, was approved in 2021.

\section*{Why This Test Is Ordered}
All patients with newly diagnosed stage IV (metastatic) or recurrent NSCLC (particularly adenocarcinoma) should undergo molecular profiling. Stage I--III NSCLC patients may also be tested (particularly in academic settings) to guide adjuvant therapy (osimertinib for EGFR-mutated resected NSCLC per the ADAURA trial). The purpose is to identify actionable driver mutations or fusions for which FDA-approved targeted therapies exist.

\section*{Primary vs Secondary Test}
This is a mandatory primary test for all patients with advanced/metastatic NSCLC adenocarcinoma per NCCN, CAP/IASLC/AMP guidelines.

\section*{How This Test Is Performed}
Tumor tissue (core needle biopsy, surgical resection, or cell block from pleural fluid) is fixed in formalin and embedded in paraffin. Sections are cut, tumor content is assessed, and DNA/RNA is extracted. Multi-gene NGS is performed. Results are typically available in 7--14 days.

\section*{What Preparation Is Needed}
No fasting or special preparation is required. Molecular profiling requires adequate tumor tissue; the ordering provider should confirm that the biopsy or resection specimen contains sufficient tumor content (usually evaluated by the pathology laboratory) or, when tissue is inadequate or unavailable, that plasma ctDNA (liquid biopsy) testing is appropriate. There is no need to hold targeted therapy before tissue or blood collection. If testing is being repeated at progression to look for resistance mechanisms, the timing of the new biopsy or ctDNA draw relative to the current therapy should be documented.

\section*{Contraindications and Precautions}
There are no contraindications to the laboratory test itself; the risks are those of the biopsy or blood draw. Interpretive cautions: (1) a tissue sample represents only the portion of tumor tested and may not capture intratumoral or metastatic heterogeneity, so a negative result does not entirely exclude an actionable alteration. (2) A negative plasma ctDNA result does not exclude a targetable alteration, particularly with low tumor burden or non-shedding tumors. (3) PD-L1 tumor proportion score thresholds differ between assays and antibody clones, and results are only reproducible within validated cutoffs (e.g., 50\%). (4) Some NGS panels detect fusions only through RNA analysis; DNA-only panels may miss ALK, ROS1, RET, or NTRK fusions. (5) Somatic tumor testing is not designed to detect germline variants; unexpected germline findings should prompt dedicated germline testing and counseling.

\section*{Risks}
The genetic test itself carries no direct risk. Risks relate to the sampling procedure: bronchoscopy or transthoracic needle biopsy may cause pneumothorax, bleeding, or infection, and a liquid biopsy blood draw carries the usual minimal venipuncture risks (slight pain, bruising, and rarely hematoma or vasovagal syncope).

\section*{What Happens After the Test}
Results are typically available in 7--14 days and are reported as the specific genomic alterations detected and the PD-L1 tumor proportion score. The results are reviewed at a multidisciplinary tumor board and guide first-line therapy (e.g., osimertinib for EGFR-sensitizing mutations, alectinib for ALK rearrangements, targeted inhibitors for other actionable drivers). If no actionable alteration is found, PD-L1 expression guides immunotherapy and chemotherapy selection. At disease progression on targeted therapy, repeat biopsy or ctDNA testing is used to identify acquired resistance mechanisms.

\section*{Understanding Results}

\subsection*{EGFR-Mutated (Sensitizing Mutation)}
EGFR exon 19 deletion or exon 21 L858R: Osimertinib (third-generation EGFR TKI) is the preferred first-line therapy (FLAURA trial; median PFS 18.9 months vs. 10.2 months for first-generation TKIs) \cite{soria2018osimertinib}. EGFR T790M: acquired resistance mutation after first-/second-generation EGFR TKI; osimertinib is active. EGFR exon 20 insertion: Amivantamab (bispecific EGFR-MET antibody) or mobocertinib for targeted therapy.

\subsection*{ALK-Rearranged (ALK-Positive)}
Alectinib, brigatinib, or lorlatinib (second-/third-generation ALK inhibitors) are preferred first-line therapy. Median PFS >34.8 months with alectinib (ALEX trial) \cite{peters2017alectinib}. Crizotinib is no longer preferred first-line.

\subsection*{ROS1-Rearranged}
Entrectinib or crizotinib are approved first-line. Lorlatinib has activity after progression on crizotinib.

\subsection*{BRAF V600E-Mutated}
Dabrafenib plus trametinib (BRAF/MEK inhibitor).

\subsection*{NTRK1/2/3 Fusions}
Larotrectinib or entrectinib (tumor-agnostic).

\subsection*{RET-Rearranged}
Selpercatinib or pralsetinib (RET-selective inhibitors).

\subsection*{MET Exon 14 Skipping}
Capmatinib or tepotinib (MET inhibitors).

\subsection*{KRAS G12C-Mutated}
Sotorasib or adagrasib (KRAS G12C-specific inhibitors) \cite{skoulidis2021sotorasib}.

\subsection*{HER2 (ERBB2)-Mutated}
Trastuzumab deruxtecan (T-DXd).

\subsection*{No Actionable Mutation (Oncogene-Negative NSCLC)}
PD-L1 expression guides immunotherapy: PD-L1 $\geq$50\% (high): pembrolizumab monotherapy. PD-L1 1--49\% (low): Pembrolizumab + platinum doublet chemotherapy. PD-L1 <1\% (negative): Platinum doublet chemotherapy + pembrolizumab (or chemotherapy + immunotherapy combination).

\section*{Normal Results}
No actionable genomic alteration detected (oncogene-negative or ``pan-negative'' NSCLC). In adenocarcinoma, approximately 60--70\% of patients have an actionable alteration. In squamous cell carcinoma, <10\% have an actionable alteration. PD-L1 IHC: reported as Tumor Proportion Score (TPS; percentage of tumor cells with membranous staining) for pembrolizumab selection.

\section*{Follow-up Tests}
PD-L1 IHC (22C3 for pembrolizumab eligibility). Liquid biopsy (ctDNA) if tissue is insufficient or if resistance mutations (EGFR T790M, emerging ALK resistance mutations) are suspected. Repeat biopsy at disease progression on targeted therapy to identify acquired resistance mechanisms. Brain MRI (CNS metastases are common in EGFR-mutated and ALK-rearranged NSCLC). Comprehensive genomic profiling for clinical trial eligibility.

\section*{References}
\bibliographystyle{unsrtnat}
\bibliography{references}

\clearpage
\section*{Regulatory Status and Authorization Date}
Regulatory status applies to the specific assay, instrument, manufacturer, and intended use rather than to the test name alone. No single approval date can be assigned to this test category. For a commercial assay, verify the current product in the relevant regulator's database: the U.S. Food and Drug Administration (FDA) clearance, approval, or authorization; India's Central Drugs Standard Control Organisation (CDSCO) licence or permission; the European Union In Vitro Diagnostic Regulation (EU IVDR) CE certificate issued through the applicable conformity-assessment route; or the United Kingdom Medicines and Healthcare products Regulatory Agency (MHRA) registration and UKCA/CE status. Laboratory-developed tests may not have an individual FDA, CDSCO, EU IVDR, or MHRA approval date.
\clearpage
